\documentclass[12pt,a4paper]{article}

\usepackage[utf8]{inputenc}
\usepackage[T1]{fontenc}
\usepackage[spanish,es-nodecimaldot,es-noquoting]{babel}
\usepackage{graphicx}
\usepackage{geometry}
\geometry{top=2.5cm,bottom=2.5cm,left=2.5cm,right=2.5cm,headheight=15pt}
\usepackage{xcolor}
\usepackage{hyperref}
\usepackage{longtable}
\usepackage{tabularx}
\usepackage{booktabs}
\usepackage{listings}
\usepackage{fancyhdr}
\usepackage{tcolorbox}
\usepackage{titlesec}
\usepackage{enumitem}
\usepackage{caption}
\usepackage{amsmath}
\usepackage{float}

% Colores institucionales y de código
\definecolor{unsaacblue}{RGB}{0, 51, 102}
\definecolor{darkblue}{RGB}{10, 40, 95}
\definecolor{accentblue}{RGB}{25, 118, 210}
\definecolor{codegreen}{RGB}{46, 125, 50}
\definecolor{codegray}{RGB}{100, 116, 139}
\definecolor{codepurple}{RGB}{124, 58, 237}
\definecolor{backcolour}{RGB}{248, 250, 252}
\definecolor{bordercolor}{RGB}{226, 232, 240}
\definecolor{calloutbg}{RGB}{240, 249, 255}
\definecolor{calloutborder}{RGB}{56, 189, 248}

% Configuración de listings para código Dart / Flutter
\lstdefinelanguage{Dart}{
  morekeywords={
    abstract,as,assert,async,await,break,case,catch,class,const,continue,
    default,deferred,do,dynamic,else,enum,export,extends,external,factory,
    false,final,finally,for,Function,get,hide,if,implements,import,in,
    interface,is,late,library,mixin,new,null,of,on,operator,part,rethrow,
    return,set,show,static,super,switch,sync,this,throw,true,try,typedef,
    var,void,while,with,yield,override,required
  },
  morecomment=[l]{//},
  morecomment=[s]{/*}{*/},
  morestring=[b]',
  morestring=[b]",
  sensitive=true
}

\lstset{
  language=Dart,
  backgroundcolor=\color{backcolour},
  commentstyle=\color{codegreen}\itshape,
  keywordstyle=\color{unsaacblue}\bfseries,
  numberstyle=\tiny\color{codegray},
  stringstyle=\color{codepurple},
  basicstyle=\ttfamily\footnotesize,
  breakatwhitespace=false,
  breaklines=true,
  captionpos=b,
  keepspaces=true,
  numbers=left,
  numbersep=7pt,
  showspaces=false,
  showstringspaces=false,
  showtabs=false,
  tabsize=2,
  frame=single,
  rulecolor=\color{bordercolor},
  xleftmargin=15pt,
  framexleftmargin=15pt
}

\hypersetup{
  colorlinks=true,
  linkcolor=unsaacblue,
  filecolor=magenta,
  urlcolor=accentblue,
  citecolor=unsaacblue
}

% Encabezados y pies de página
\pagestyle{fancy}
\fancyhf{}
\fancyhead[L]{\small\color{codegray}UNSAAC - Desarrollo de Software II}
\fancyhead[R]{\small\color{codegray}Práctica 06: Diseño Responsivo y Accesible en Flutter}
\fancyfoot[C]{\thepage}
\renewcommand{\headrulewidth}{0.4pt}
\renewcommand{\footrulewidth}{0.4pt}

\begin{document}

% ==============================================================================
% CARÁTULA OFICIAL INSTITUCIONAL UNSAAC
% ==============================================================================
\begin{titlepage}
    \centering
    {\bfseries\large UNIVERSIDAD NACIONAL DE SAN ANTONIO ABAD DEL CUSCO\par}
    \vspace{0.2cm}
    {\bfseries\normalsize FACULTAD DE INGENIERÍA ELÉCTRICA, ELECTRÓNICA, INFORMÁTICA Y MECÁNICA\par}
    \vspace{0.2cm}
    {\bfseries\normalsize ESCUELA PROFESIONAL DE INGENIERÍA INFORMÁTICA Y DE SISTEMAS\par}
    \vspace{0.8cm}
    \includegraphics[width=0.35\textwidth]{unsaac.jpg}\par
    \vspace{0.8cm}
    {\Large\bfseries ASIGNATURA: DESARROLLO DE SOFTWARE II\par}
    \vspace{0.4cm}
    {\large\bfseries TRABAJO: PRÁCTICA 06: DISEÑO RESPONSIVO Y ACCESIBLE EN FLUTTER\par}
    \vspace{0.4cm}
    {\large\bfseries IMPLEMENTACIÓN INTEGRAL DE SISTEMA RESPONSIVO, ADAPTABILIDAD MULTIPLATAFORMA, NORMAS WCAG AA Y ACCESIBILIDAD SEMÁNTICA\par}
    \vspace{1.2cm}
    \begin{flushleft}
    \textbf{DOCENTE RESPONSABLE:} \\
    Ing. Hans Harley Ccacyahuillca Bejar\\[0.5cm]
    \textbf{AUTOR / DESARROLLADOR:} \\
    SUPA CUSIPAUCAR, Yeferson - 220553\\[0.5cm]
    \textbf{SEMESTRE ACADÉMICO:} \\
    2026-II (Décimo Semestre)\\[0.5cm]
    \textbf{REPOSITORIO OFICIAL:} \\
    \url{https://github.com/YefersonSupa/Desarrollo_SoftwareII}
    \end{flushleft}
    \vfill
    {\bfseries\large Cusco - Perú\par}
    {\bfseries 2026\par}
\end{titlepage}

\tableofcontents
\newpage

% ==============================================================================
% 1. INFORMACIÓN GENERAL Y FICHA TÉCNICA
% ==============================================================================
\section{Información General del Proyecto}

El presente informe técnico documenta el diseño, la formulación arquitectónica, la implementación de código fuente y la validación automatizada correspondiente a la \textbf{Práctica de Laboratorio N$^\circ$ 06: Diseño Responsivo y Accesible en Flutter} de la cátedra de \textbf{Desarrollo de Software II}.

El desarrollo ha sido articulado a través de un aplicativo modular denominado \texttt{responsive\_app}, concebido para materializar de manera exhaustiva todos y cada uno de los lineamientos dictados en la guía de cátedra:
\begin{enumerate}[leftmargin=*]
    \item \textbf{Mecanismos de Layout Responsivo:} Utilización rigurosa de \texttt{LayoutBuilder}, \texttt{MediaQuery.sizeOf(context)}, \texttt{FractionallySizedBox}, y combinatorias proporcionales de \texttt{Expanded} y \texttt{Flexible} con factores \texttt{flex}.
    \item \textbf{Breakpoints y Diseño Adaptativo:} Definición centralizada de umbrales ergonómicos en la clase \texttt{Breakpoints} (móvil $<480\text{ px}$, tablet $480-1024\text{ px}$, escritorio $\ge 1024\text{ px}$, ultrapanorámico $\ge 1440\text{ px}$), grilla adaptable con \texttt{SliverGridDelegateWithMaxCrossAxisExtent} y navegación bifurcada mediante \texttt{NavigationRail} en pantallas amplias frente a \texttt{NavigationBar} en dispositivos móviles.
    \item \textbf{Directrices de Accesibilidad Visual y Motriz (WCAG AA):} Cumplimiento de ratios de contraste cromático $\ge 4.5:1$ en texto estándar, áreas mínimas de contacto táctil (\textit{touch targets}) $\ge 48\times 48\text{ px}$, escalamiento de texto libre de bloqueos forzados de \texttt{textScaleFactor} respetando las preferencias del sistema operativo, y navegación por teclado guiada mediante \texttt{FocusTraversalOrder}.
    \item \textbf{Semántica para Tecnologías Asistivas:} Configuración explícita del árbol semántico para lectores de pantalla (\textit{TalkBack} en Android, \textit{VoiceOver} en iOS y lectores de escritorio) mediante widgets \texttt{Semantics}, exclusión de ornamentos visuales con \texttt{ExcludeSemantics}, fusión de controles compuestos mediante \texttt{MergeSemantics}, anuncios vocales dinámicos con \texttt{SemanticsService.announce} e inspección en vivo con \texttt{SemanticsDebugger}.
    \item \textbf{Lista de Verificación de Lanzamiento (\textit{Checklist}):} Módulo interactivo con el inventario completo de los \textbf{20 criterios oficiales} tabulados en la guía de práctica, clasificados en cuatro ejes: Diseño Responsivo, Tipografía, Accesibilidad y Pruebas.
\end{enumerate}

\begin{tcolorbox}[colback=calloutbg,colframe=calloutborder,title=\textbf{Ficha Técnica de la Práctica 06}]
\textbf{Institución:} Universidad Nacional de San Antonio Abad del Cusco (UNSAAC)\\
\textbf{Facultad:} Ingeniería Eléctrica, Electrónica, Informática y Mecánica\\
\textbf{Escuela Profesional:} Ingeniería Informática y de Sistemas\\
\textbf{Asignatura:} Desarrollo de Software II\\
\textbf{Docente Responsable:} Ing. Hans Harley Ccacyahuillca Bejar\\
\textbf{Estudiante:} SUPA CUSIPAUCAR, Yeferson (Código: 220553)\\
\textbf{Tecnologías Clave:} Flutter 3.47.4, Dart 3.13.3, Material Design 3, WCAG 2.1 AA\\
\textbf{Ubicación del Código Fuente:} \texttt{/lab\_06/responsive_app}\\
\textbf{Repositorio en GitHub:} \url{https://github.com/YefersonSupa/Desarrollo_SoftwareII}
\end{tcolorbox}

% ==============================================================================
% 2. FUNDAMENTOS TEÓRICOS DE RESPONSIVIDAD Y ACCESIBILIDAD
% ==============================================================================
\section{Fundamentos Teóricos de Responsividad y Accesibilidad en Flutter}

\subsection{El Algoritmo de Diseño (\textit{Layout}) de Flutter y sus Restricciones}
A diferencia de los modelos basados en hojas de estilo en cascada (CSS) donde los elementos calculan sus dimensiones a partir de reglas implícitas o flujos de caja, Flutter opera bajo un axioma estricto y unidireccional:
\begin{quote}
\textit{Constraints go down. Sizes go up. Parent sets position.}
\end{quote}
Las restricciones dimensionales (\texttt{BoxConstraints}: valores mínimos y máximos de ancho y alto) descienden desde el nodo padre hacia el hijo. El nodo hijo, conociendo tales límites, determina su tamaño final (\texttt{Size}) y lo devuelve a su progenitor. Finalmente, el padre decide las coordenadas relativas $(x, y)$ donde ubicar al hijo.

\subsection{Diferenciación Técnica: MediaQuery frente a LayoutBuilder}
Uno de los errores más comunes en el desarrollo responsivo consiste en acoplar el comportamiento interno de los componentes al tamaño de la pantalla global mediante \texttt{MediaQuery}:
\begin{itemize}[leftmargin=*]
    \item \textbf{MediaQuery.sizeOf(context):} Proporciona la resolución total de la ventana gráfica (\textit{viewport}) de la aplicación. Es adecuado para decisiones arquitecturales de alto nivel (como conmutar entre una barra de navegación inferior \texttt{NavigationBar} y un riel lateral \texttt{NavigationRail}). Sin embargo, suscribir widgets secundarios a \texttt{MediaQuery} provoca reconstrucciones innecesarias cada vez que cambia cualquier métrica periférica del dispositivo.
    \item \textbf{LayoutBuilder:} Entrega al constructor las \texttt{BoxConstraints} exactas delegadas por el widget padre inmediato. Esto permite diseñar componentes altamente reutilizables y desacoplados: una misma tarjeta puede adaptarse de vista vertical a horizontal según el ancho del contenedor en el que se inserte (por ejemplo, dentro de un panel lateral o en el área de contenido principal), independientemente de las dimensiones globales de la pantalla.
\end{itemize}

\subsection{Widgets Clave de Adaptabilidad Proporcional}
\begin{itemize}[leftmargin=*]
    \item \textbf{FractionallySizedBox:} Permite que un widget ocupe un porcentaje explícito ($0.0$ a $1.0$) del espacio disponible en su padre, facilitando la diagramación fluida sin necesidad de cálculos matemáticos manuales dependientes de constantes en píxeles lógicos.
    \item \textbf{Expanded y Flexible:} Dentro de una fila (\texttt{Row}) o columna (\texttt{Column}), distribuyen el espacio remanente mediante el parámetro escalar \texttt{flex}. \texttt{Expanded} obliga al hijo a llenar exactamente todo el espacio asignado por su factor de escala (\texttt{fit: FlexFit.tight}), mientras que \texttt{Flexible} le otorga un techo máximo sin obligarlo a expandirse artificialmente (\texttt{FlexFit.loose}).
    \item \textbf{SliverGridDelegateWithMaxCrossAxisExtent:} Permite que una grilla calcule automáticamente el número óptimo de columnas en tiempo de ejecución. Al fijar el ancho máximo que un ítem puede abarcar (por ejemplo, $300\text{ px}$), la grilla añade o retira columnas dinámicamente conforme la ventana cambia de tamaño.
\end{itemize}

\subsection{Estándares de Accesibilidad Digital (WCAG 2.1 AA)}
La accesibilidad no es un añadido estético, sino un requisito de ingeniería indispensable para garantizar la universalidad del software:
\begin{enumerate}[leftmargin=*]
    \item \textbf{Contraste Cromático:} La pauta WCAG 2.1 AA exige un ratio de contraste mínimo de $4.5:1$ entre el color del texto y su fondo para texto regular (menor a $18\text{ pt}$ o $14\text{ pt}$ negrita), y de $3.0:1$ para texto de gran tamaño o componentes de interfaz esenciales.
    \item \textbf{Áreas Táctiles Mínimas (\textit{Touch Targets}):} Las directrices de ergonomía de Google Material Design y Apple Human Interface Guidelines dictan que todo control interactivo debe ofrecer un objetivo táctil de al menos $48\times 48\text{ px}$ lógicos ($44\times 44\text{ pt}$ en iOS), evitando fallos en usuarios con destreza motriz reducida.
    \item \textbf{Escalamiento Dinámico de Texto:} El sistema debe permitir al usuario ajustar el tamaño de fuente sin que la aplicación trunque textos, rompa contenedores ni fuerce un \texttt{textScaleFactor} estático.
    \item \textbf{Árbol Semántico y Tecnologías Asistivas:} Los lectores de pantalla navegan mediante una representación estructural paralela al árbol visual (\textit{Semantics Tree}). Flutter provee widgets como \texttt{Semantics} (para dotar de etiquetas, pistas de uso y roles), \texttt{ExcludeSemantics} (para ocultar iconos puramente estéticos) y \texttt{MergeSemantics} (para unificar controles interactivos compuestos y evitar redundancia auditiva).
\end{enumerate}

% ==============================================================================
% 3. ARQUITECTURA DE LA APLICACIÓN RESPONSIVE_APP
% ==============================================================================
\section{Arquitectura y Organización del Proyecto}

El proyecto ha sido estructurado bajo el principio de separación de responsabilidades y modularidad en carpetas limpias dentro del paquete Flutter \texttt{lab\_06/responsive\_app}:

\begin{figure}[H]
\centering
\begin{tcolorbox}[colback=backcolour,colframe=unsaacblue,title=\textbf{Organización del Código Fuente - lab\_06/responsive\_app}]
\small
\begin{verbatim}
lab_06/
├── Guia06-flutter_responsivo_accesible.docx.pdf # Guía oficial de cátedra
├── unsaac.jpg                                   # Escudo institucional UNSAAC
├── informe_laboratorio_06.tex                   # Informe formal en LaTeX
├── README.md                                    # Guía técnica de ejecución
├── ej1.png ... ej5.png                          # Capturas de evidencia generadas
└── responsive_app/                              # Proyecto Flutter ejecutable
    ├── lib/
    │   ├── main.dart                            # Punto de entrada y Scaffold adaptativo
    │   ├── utils/
    │   │   └── breakpoints.dart                 # Centralización de breakpoints y grids
    │   ├── models/
    │   │   ├── catalog_item.dart                # Modelo de datos y dataset semántico
    │   │   └── checklist_item.dart              # Modelo y 20 criterios de lanzamiento
    │   └── screens/
    │       ├── responsive_layout_screen.dart    # Módulo 1: LayoutBuilder y MediaQuery
    │       ├── adaptive_grid_screen.dart        # Módulo 2: Grillas y Breakpoints
    │       ├── accessibility_screen.dart        # Módulo 3: WCAG AA, Contraste y Texto
    │       ├── semantics_screen.dart            # Módulo 4: Semantics y Lectores
    │       └── checklist_screen.dart            # Módulo 5: Checklist interactivo
    └── test/
        ├── responsive_accessible_test.dart      # 9 pruebas unitarias y de widgets
        └── capture_screenshots_test.dart        # Generador automatizado de capturas
\end{verbatim}
\end{tcolorbox}
\caption{Estructura de directorios y segregación de módulos del proyecto.}
\label{fig:estructura}
\end{figure}

% ==============================================================================
% 4. DESARROLLO E IMPLEMENTACIÓN DE MÓDULOS
% ==============================================================================
\section{Desarrollo e Implementación de los Módulos de la Guía}

\subsection{Módulo 1: Layout Responsivo (MediaQuery y LayoutBuilder)}
El primer módulo ilustra cómo gestionar el espacio gráfico mediante widgets adaptables nativos. Implementa:
\begin{itemize}[leftmargin=*]
    \item Un panel de telemetría en tiempo real que monitoriza el ancho, alto y orientación del \textit{viewport} usando \texttt{MediaQuery.sizeOf(context)}.
    \item Un contenedor adaptable gestionado con \texttt{LayoutBuilder} que commuta de disposición vertical (\texttt{Column}) a horizontal (\texttt{Row}) cuando el ancho disponible supera el umbral de $600\text{ px}$.
    \item Un control deslizante que ajusta dinámicamente un \texttt{FractionallySizedBox} entre $20\%$ y $100\%$ del ancho asignado.
    \item Un comparador proporcional que conmuta entre proporciones \texttt{Expanded} equilibradas ($1:1$) y asimétricas ($2:1$).
\end{itemize}

\begin{figure}[H]
\centering
\includegraphics[width=0.65\textwidth]{ej1.png}
\caption{Evidencia Módulo 1: Layout Responsivo interactivo con LayoutBuilder y MediaQuery.}
\label{fig:ej1}
\end{figure}

\begin{lstlisting}[caption={Fragmento de conmutación adaptable con LayoutBuilder}]
LayoutBuilder(
  builder: (context, constraints) {
    final isWide = constraints.maxWidth > 600;
    if (isWide) {
      return Row(
        crossAxisAlignment: CrossAxisAlignment.start,
        children: [
          Expanded(child: _buildInfoCard(context, "Panel Primario", ...)),
          const SizedBox(width: 16),
          Expanded(child: _buildInfoCard(context, "Panel Secundario", ...)),
        ],
      );
    }
    return Column(
      children: [
        _buildInfoCard(context, "Panel Primario", ...),
        const SizedBox(height: 12),
        _buildInfoCard(context, "Panel Secundario", ...),
      ],
    );
  },
)
\end{lstlisting}

\subsection{Módulo 2: Grilla Adaptativa y Breakpoints}
El segundo módulo aplica la clase centralizada \texttt{Breakpoints}:
\begin{itemize}[leftmargin=*]
    \item \textbf{Breakpoints Centralizados:} Móvil ($<480\text{ px}$), Tablet ($480-1024\text{ px}$), Desktop ($\ge 1024\text{ px}$) y Ultra-ancho ($\ge 1440\text{ px}$).
    \item \textbf{Grilla Adaptativa Dinámica:} Utiliza \texttt{SliverGridDelegateWithMaxCrossAxisExtent} con un control interactivo que permite al evaluador alterar el valor de \texttt{maxCrossAxisExtent} entre $160\text{ px}$ y $340\text{ px}$, observando en vivo cómo las columnas se reorganizan automáticamente.
    \item \textbf{Dataset Semántico de Servicios Urbanos:} 8 tarjetas de servicios ciudadanos (\textit{Alumbrado Público, Limpieza Urbana, Mantenimiento de Parques, Seguridad Vecinal, etc.}) enriquecidas con iconos, estados de urgencia, valoraciones numéricas y resúmenes descriptivos.
\end{itemize}

\begin{figure}[H]
\centering
\includegraphics[width=0.65\textwidth]{ej2.png}
\caption{Evidencia Módulo 2: Grilla Adaptativa con Breakpoints dinámicos y dataset de servicios.}
\label{fig:ej2}
\end{figure}

\subsection{Módulo 3: Accesibilidad, Contraste y Escalamiento de Texto}
El tercer módulo aborda la satisfacción de las pautas WCAG 2.1 AA mediante tres laboratorios interactivos integrados:
\begin{itemize}[leftmargin=*]
    \item \textbf{Simulador de Escalamiento de Fuentes:} Permite modificar el multiplicador de escala tipográfica desde $0.85\times$ hasta $1.8\times$ en pasos discretos, demostrando que los componentes no desbordan ni truncan contenido.
    \item \textbf{Auditor de Objetivos Táctiles (\textit{Touch Targets}):} Muestra una comparación directa entre un botón defectuoso de $24\times 24\text{ px}$ (inadecuado para destreza motriz reducida) y un botón accesible que cumple holgadamente la cota de $48\times 48\text{ px}$.
    \item \textbf{Modo Alto Contraste (WCAG AA):} Conmutador que transforma el tema visual aplicando negro azabache de fondo (\texttt{\#000000}) y texto amarillo de alta luminosidad (\texttt{\#FFD700}), alcanzando un ratio de contraste superior a $15:1$.
    \item \textbf{Navegación por Teclado:} Orden de foco explícito implementado mediante \texttt{FocusTraversalGroup} y \texttt{NumericFocusOrder}.
\end{itemize}

\begin{figure}[H]
\centering
\includegraphics[width=0.65\textwidth]{ej3.png}
\caption{Evidencia Módulo 3: Auditoría de Accesibilidad, Touch Targets y Escalamiento.}
\label{fig:ej3}
\end{figure}

\subsection{Módulo 4: Semántica para Lectores de Pantalla}
El cuarto módulo demuestra la interacción con motores asistivos (\textit{Screen Readers}):
\begin{itemize}[leftmargin=*]
    \item \textbf{Semantics:} Asignación de etiquetas personalizadas (\texttt{label}), pistas de interacción (\texttt{hint}) e indicación explícita de botón interactivo (\texttt{button: true}) sobre un avatar de perfil.
    \item \textbf{ExcludeSemantics:} Ocultamiento de iconos decorativos y estrellas de valoración para evitar que el lector de pantalla enuncie redundantemente cada elemento gráfico no sustancial.
    \item \textbf{MergeSemantics:} Agrupación del checkbox de aceptación de términos con su texto descriptivo, permitiendo que el usuario con discapacidad visual interactúe con el componente como una sola entidad cohesiva.
    \item \textbf{SemanticsService.announce y liveRegion:} Disparo de mensajes vocales prioritarios que notifican eventos asíncronos y contenedor con \texttt{liveRegion: true} para lectura automática al mutar su contenido.
    \item \textbf{SemanticsDebugger:} Interruptor en vivo que activa el renderizado de la superposición diagnóstica de Flutter, visualizando los rectángulos del árbol semántico sobre la interfaz gráfica.
\end{itemize}

\begin{figure}[H]
\centering
\includegraphics[width=0.65\textwidth]{ej4.png}
\caption{Evidencia Módulo 4: Semántica avanzada para lectores de pantalla y SemanticsDebugger.}
\label{fig:ej4}
\end{figure}

\subsection{Módulo 5: Checklist Interactivo de Lanzamiento}
El quinto módulo materializa los 20 criterios oficiales del checklist de lanzamiento estipulados en la página 8 de la Guía 06:
\begin{itemize}[leftmargin=*]
    \item Barra de progreso porcentual y conteo de criterios satisfechos en tiempo real.
    \item Filtros por categoría: \textit{Todos, Responsivo, Tipografía, Accesibilidad, Pruebas}.
    \item Botón de marcado y desmarcado global con persistencia de estado durante la sesión interactiva.
\end{itemize}

\begin{figure}[H]
\centering
\includegraphics[width=0.65\textwidth]{ej5.png}
\caption{Evidencia Módulo 5: Lista de verificación (\textit{Checklist}) oficial de 20 criterios.}
\label{fig:ej5}
\end{figure}

% ==============================================================================
% 5. MATRIZ DE VERIFICACIÓN DEL CHECKLIST DE LANZAMIENTO
% ==============================================================================
\section{Matriz de Verificación del Checklist Oficial de Lanzamiento}

A continuación, se documenta la matriz completa de los 20 criterios oficiales exigidos para el pase a producción de interfaces en Flutter, implementados y verificados en la aplicación \texttt{responsive\_app}. Cumpliendo con los estándares institucionales, esta tabla se formula mediante el entorno \texttt{longtable} con control dimensional estricto para evitar desbordes de margen:

\begin{small}
\begin{longtable}{p{0.8cm} p{2.5cm} p{5.5cm} p{3.2cm} p{1.8cm}}
\caption{Matriz oficial de los 20 criterios del Checklist de Lanzamiento (Guía 06).} \label{tab:checklist_lanzamiento} \\
\toprule
\textbf{N$^\circ$} & \textbf{Categoría} & \textbf{Criterio de Verificación} & \textbf{Widget / Mecanismo} & \textbf{Estado} \\
\midrule
\endfirsthead

\multicolumn{5}{c}%
{{\bfseries \tablename\ \thetable{} -- Continuación de la página anterior}} \\
\toprule
\textbf{N$^\circ$} & \textbf{Categoría} & \textbf{Criterio de Verificación} & \textbf{Widget / Mecanismo} & \textbf{Estado} \\
\midrule
\endhead

\midrule
\multicolumn{5}{r}{{Continúa en la siguiente página \dots}} \\
\bottomrule
\endfoot

\bottomrule
\endlastfoot

01 & Responsivo & Sin desbordamientos (\textit{overflow errors}) en pantallas estándar & \texttt{LayoutBuilder}, \texttt{Expanded} & Conforme \\
02 & Responsivo & Disposición adaptada en pantallas anchas (tablet/desktop) & \texttt{NavigationRail}, \texttt{Row} & Conforme \\
03 & Responsivo & Navegación adaptada según factor de forma & \texttt{NavigationBar} / \texttt{Rail} & Conforme \\
04 & Responsivo & Grillas con columnas dinámicas según ancho & \texttt{MaxCrossAxisExtent} & Conforme \\
05 & Responsivo & Elementos interactivos utilizables en orientación horizontal & \texttt{SingleChildScrollView} & Conforme \\
06 & Tipografía & No se bloquea el escalamiento del sistema & \texttt{TextScaler} dinámico & Conforme \\
07 & Tipografía & Textos no se truncan con escalamiento hasta $1.5\times$ & \texttt{Flexible}, \texttt{Wrap} & Conforme \\
08 & Tipografía & Se usa \texttt{TextScaler.linear} para fijar límites controlados & \texttt{TextScaler.clamp} & Conforme \\
09 & Tipografía & Fuentes legibles con pesos tipográficos diferenciados & \texttt{FontWeight} (M3) & Conforme \\
10 & Tipografía & Altura de línea (\texttt{height}) configurada para evitar colisiones & \texttt{TextStyle.height: 1.3} & Conforme \\
11 & Accesibilidad & Contraste texto/fondo cumple WCAG AA ($4.5:1$ texto normal) & \texttt{ColorScheme} optimizado & Conforme \\
12 & Accesibilidad & Contraste texto/fondo cumple WCAG AA ($3.0:1$ texto grande) & \texttt{ThemeData.highContrast} & Conforme \\
13 & Accesibilidad & Objetivos táctiles mínimos de $48\times 48\text{ px}$ lógicos & Botones con \texttt{kMinInteractive} & Conforme \\
14 & Accesibilidad & Imágenes significativas cuentan con etiqueta semántica & \texttt{Semantics.label} & Conforme \\
15 & Accesibilidad & Iconos decorativos excluidos del lector de pantalla & \texttt{ExcludeSemantics} & Conforme \\
16 & Accesibilidad & Elementos complejos fusionados con \texttt{MergeSemantics} & \texttt{MergeSemantics} en tiles & Conforme \\
17 & Accesibilidad & Cambios de estado anunciados con \texttt{SemanticsService} & \texttt{SemanticsService.announce} & Conforme \\
18 & Accesibilidad & Orden de foco lógico para navegación por teclado & \texttt{FocusTraversalOrder} & Conforme \\
19 & Pruebas & Probado con lector de pantalla (TalkBack / VoiceOver) & \texttt{SemanticsDebugger} & Conforme \\
20 & Pruebas & Probado con suite de tests automatizados & \texttt{flutter test} (9 pruebas) & Conforme \\

\end{longtable}
\end{small}

% ==============================================================================
% 6. SUITE DE PRUEBAS AUTOMATIZADAS
% ==============================================================================
\section{Suite de Pruebas Automatizadas y Verificación}

Para asegurar la solidez del sistema ante regresiones y verificar los criterios técnicos sin ambigüedad, se formuló una suite completa de pruebas unitarias y de widgets en el archivo \texttt{test/responsive\_accessible\_test.dart}. La batería consta de 9 casos de prueba exhaustivos que validan:
\begin{enumerate}[leftmargin=*]
    \item \textbf{Lógica de Breakpoints:} Cálculo matemático de columnas para resoluciones móviles ($360\text{ px}$), tablets intermedias ($720\text{ px}$) y escritorios ($1200\text{ px}$).
    \item \textbf{Integridad de Modelos:} Verificación de que el dataset de servicios contenga 8 ítems íntegros y que la lista de chequeo incluya con exactitud los 20 criterios oficiales.
    \item \textbf{Renderizado Inicial y Navegación Móvil:} Comprobación de la barra de navegación inferior \texttt{NavigationBar} con sus 5 pestañas activas.
    \item \textbf{Adaptabilidad a Escritorio:} Simulación de pantalla de $1200\times 800\text{ px}$ y aserción de la conmutación a \texttt{NavigationRail}.
    \item \textbf{Inspección de Grilla Adaptativa:} Verificación de la presencia de las tarjetas de servicios urbanos y renderizado sin desbordamientos de renderizado (\textit{zero RenderFlex overflow}).
    \item \textbf{Auditoría de Touch Targets:} Medición precisa de la caja de colisión del botón accesible, asegurando dimensiones mínimas $\ge 48\times 48\text{ px}$.
    \item \textbf{Verificación Semántica y MergeSemantics:} Comprobación de que el árbol semántico del widget fusionado reporte el nodo compuesto correctamente.
    \item \textbf{Interactividad del Checklist:} Prueba de alternancia de estado al pulsar sobre el primer criterio de verificación.
\end{enumerate}

\begin{tcolorbox}[colback=backcolour,colframe=codegreen,title=\textbf{Resultado de Ejecución de Pruebas Automatizadas}]
\ttfamily\small
00:00 +0: Breakpoints: Calculo de columnas dinamicas para grilla adaptativa\\
00:00 +1: CatalogItem: Dataset de muestra contiene 8 elementos con semantica completa\\
00:00 +2: ChecklistItem: Los 20 criterios oficiales de la Guia 06 estan presentes\\
00:00 +3: 1. Renderizado inicial de la aplicacion y navegacion movil\\
00:00 +4: 2. Adaptabilidad Desktop: Renderiza NavigationRail en pantallas anchas\\
00:00 +5: 3. Navegacion a Adaptive Grid y presencia de tarjetas\\
00:01 +6: 4. Verificacion de Touch Target optimo (48x48 px)\\
00:01 +7: 5. Verificacion de Semantica y MergeSemantics\\
00:01 +8: 6. Verificacion de Checklist interactivo y toggle de items\\
\textbf{00:01 +9: All tests passed!}
\end{tcolorbox}

% ==============================================================================
% 7. CONCLUSIONES Y LECCIONES APRENDIDAS
% ==============================================================================
\section{Conclusiones y Lecciones Aprendidas}

\begin{enumerate}[leftmargin=*]
    \item \textbf{Superioridad del desacoplamiento con LayoutBuilder:} El empleo de \texttt{LayoutBuilder} frente a \texttt{MediaQuery} permite desarrollar micro-componentes verdaderamente reutilizables, los cuales basan su lógica de presentación en las restricciones del contenedor inmediato y no en la resolución global del dispositivo.
    \item \textbf{La accesibilidad como pilar de arquitectura:} Diseñar interfaces accesibles en Flutter no supone un sobrecoste si se planifica desde el inicio. El uso de temas accesibles con contrastes superiores a $4.5:1$, el respeto al escalamiento del sistema y la configuración semántica mediante \texttt{Semantics} y \texttt{MergeSemantics} elevan la calidad del producto final y habilitan el soporte multiplataforma genuino.
    \item \textbf{Garantía de calidad con Touch Targets y Breakpoints Centralizados:} Centralizar los puntos de ruptura en una clase estática (\texttt{Breakpoints}) elimina valores fijos dispersos por el código fuente (\textit{magic numbers}). Asimismo, asegurar un tamaño mínimo de interacción de $48\times 48\text{ px}$ previene errores de digitación en dispositivos táctiles.
    \item \textbf{Efectividad del Checklist de 20 Criterios:} La integración del checklist interactivo dentro de la propia aplicación proporciona a los equipos de ingeniería y control de calidad (QA) un instrumento tangible y auditable para verificar la madurez del software antes del despliegue productivo.
\end{enumerate}

% ==============================================================================
% 8. REFERENCIAS BIBLIOGRÁFICAS
% ==============================================================================
\section{Referencias Bibliográficas}

\begin{enumerate}[leftmargin=*]
    \item Flutter Documentation. (2026). \textit{Creating responsive and adaptive apps}. Google LLC. Disponible en: \url{https://docs.flutter.dev/ui/layout/responsive}
    \item Flutter Documentation. (2026). \textit{Accessibility in Flutter}. Google LLC. Disponible en: \url{https://docs.flutter.dev/ui/accessibility-and-localization/accessibility}
    \item W3C. (2018). \textit{Web Content Accessibility Guidelines (WCAG) 2.1 - Contrast (Minimum) Level AA}. World Wide Web Consortium. Disponible en: \url{https://www.w3.org/WAI/WCAG21/Understanding/contrast-minimum.html}
    \item Google Material Design 3. (2026). \textit{Layout and Responsive Grid System}. Google Inc. Disponible en: \url{https://m3.material.io/foundations/layout}
\end{enumerate}

\end{document}
